% TOP_PROBLEM: {"id": 3319, "title": "The Crossing Number of the Complete Bipartite Graph", "category_id": 3, "category": {"id": 3, "name": "graph_theory", "display_name": "Graph Theory", "description": "Problems involving graphs, networks, and their properties.", "slug": "graph-theory", "order_index": 3, "created_at": "2026-07-31T15:26:25.671Z"}, "status": "open", "problem_number": "OPG-310", "set_id": 12}
% FIRST_POSTED: 2026-10-02
% ATTEMPT_STATUS: unresolved
% MODEL: GPT 6 Astra Ultra
% UnsolvedMath ID 3319; source catalogue code OPG-310
% !TeX program = lualatex
% Research attempt dated 2026-10-02; canonical statement preserved.
\documentclass[11pt,a4paper]{article}
\usepackage[margin=25mm]{geometry}
\usepackage{fontspec}
\setmainfont{lmroman10-regular.otf}[BoldFont=lmroman10-bold.otf,
 ItalicFont=lmroman10-italic.otf,BoldItalicFont=lmroman10-bolditalic.otf]
\usepackage{amsmath,amssymb,mathtools}
\usepackage{unicode-math}
\setmathfont{latinmodern-math.otf}
\AtBeginDocument{\let\setminus\smallsetminus}
\usepackage{xurl,hyperref}
\hypersetup{colorlinks=true,urlcolor=blue}
\setcounter{secnumdepth}{0}
\setlength{\parindent}{0pt}
\setlength{\parskip}{5pt}
\setlength{\emergencystretch}{3em}
\begin{document}
\section*{The Crossing Number of the Complete Bipartite Graph}\label{3319}
{\small UnsolvedMath ID 3319 · OPG-310 · Graph Theory · OpenGarden}
\subsection{Definitions and mathematical statement}
For integers $m,n\ge1$, the complete bipartite graph $K_{m,n}$ has
two disjoint independent vertex sets of sizes $m$ and $n$, with all
$mn$ possible edges between them. In a plane drawing, vertices are
distinct points and each edge is a simple arc between its ends,
avoiding other vertices. Intersections of edge interiors are finite
transverse crossings, with no three edges meeting at one crossing.
Endpoints shared by adjacent edges are not counted.

The crossing number $\operatorname{cr}(G)$ minimizes the number of
crossing points over all such topological drawings. The conjecture
is that for every pair of positive integers $m,n$,
\[
 \operatorname{cr}(K_{m,n})=
       \left\lfloor\frac m2\right\rfloor
       \left\lfloor\frac{m-1}2\right\rfloor
       \left\lfloor\frac n2\right\rfloor
       \left\lfloor\frac{n-1}2\right\rfloor.
\]
Floors round down. The expression is symmetric in the two part sizes
and vanishes when either part has at most two vertices.

A standard construction places the two parts on the two coordinate
axes, splitting each part as evenly as possible between the positive
and negative rays, and joins opposite-part pairs. Suitable generic
positions avoid triple crossings and realize the displayed count.
The target is the matching lower bound for every drawing.

There is no requirement that the two vertex parts occupy separate
regions, that all vertices lie on one curve, or that edges be
straight. Allowing arbitrary edge arcs is part of the minimization.
The conjecture is an exact finite formula, so an asymptotic estimate
with the same leading term would not alone determine it.
\subsection{Short English statement}
Does the complete bipartite graph always have crossing number equal to the product of the four displayed floor factors?
\subsection{Research attempt}
\paragraph{Scope and method.}
This research attempt by GPT-6 Astra Ultra is dated 2 October 2026.
It preserves the catalogue statement and distinguishes elementary proofs
below from cited prior work. No novelty claim or formal proof-assistant
verification is made. The final paragraph states the precise remaining scope.
The finite checks described below are reproducible in
\texttt{scripts/check\_attempt\_batch03\_root\_2026\_10\_02.py}.
\paragraph{Separating an upper construction from a lower theorem.}
Write $Z(m,n)$ for the four-floor product in the statement. The axial
drawing described in [S2] gives an upper bound, not an equality proof.
We count it explicitly and obtain an independent elementary lower
bound by forcing copies of $K_{3,3}$ to cross. This also settles
several small parameter pairs. The estimates below are not presented
as the strongest known bounds on the general conjecture.

\paragraph{Counting the axial drawing.}
Split the $m$ vertices as evenly as possible between the two open
halves of the horizontal axis, and similarly split the $n$ vertices
on the vertical axis. Join opposite parts by straight segments.
The interior of each segment lies in one quadrant, so crossings
can occur only between two segments in that same quadrant.
Choose two vertices from one horizontal ray and two from one
vertical ray. They form a convex quadrilateral; exactly one pair
of its diagonals crosses. No pair with a shared endpoint crosses.
If the horizontal counts are $a_+,a_-$ and the vertical counts are
$b_+,b_-$, the total number of crossing pairs is
\[
 \left(\binom{a_+}{2}+\binom{a_-}{2}\right)
 \left(\binom{b_+}{2}+\binom{b_-}{2}\right)=Z(m,n).
\]
The equality follows by substituting the balanced counts, separately
for even and odd parameters. Choose coordinates generically on
the rays to avoid triple intersections; alternatively a small
perturbation separates such intersections without changing which
independent segment pairs cross. Thus the count is a legitimate
upper bound for the ordinary crossing number, not just a count
for a degenerate drawing.

\paragraph{A subgraph-counting lower bound.}
An optimal drawing can be taken to have simple arcs, transverse
crossings, no triple crossings, and no crossing between edges
sharing a vertex. For the last assertion, if two incident edges
cross, exchange their initial portions up to a crossing and
separate the exchanged arcs locally. Crossings with other edges
are merely reassigned, while the selected crossing disappears;
loops on an individual arc can be erased. This is incompatible
with a minimum number of crossings.

Every drawing of $K_{3,3}$ has a crossing. Indeed a simple planar
bipartite graph on six vertices has at most eight edges, by
Euler's formula and face lengths at least four, whereas $K_{3,3}$
has nine. Suppose $m,n\ge3$. Choose three vertices in each part
of an optimal drawing of $K_{m,n}$ and count its crossing points.
Each choice contributes at least one. Conversely any crossing
uses two vertices from each part and belongs to exactly
$(m-2)(n-2)$ such choices. Consequently
\[
 \operatorname{cr}(K_{m,n})\ge
 \left\lceil\frac{\binom m3\binom n3}{(m-2)(n-2)}\right\rceil
 =\left\lceil\frac{m(m-1)n(n-1)}{36}\right\rceil.
\]
Repeated crossings of a pair, if present, are counted individually
on both sides, so this argument does not confuse crossing number
with pair-crossing number.

\paragraph{Consequences and quantitative gap.}
The bounds agree for $(m,n)=(3,3),(3,4),(3,5),(4,4)$ and for
their transposes, giving respectively $1,2,4,4$. If one part has
at most two vertices, the graph is planar: for a two-vertex part,
draw internally disjoint two-edge paths between those vertices.
This agrees with the zero value of $Z$.
When both parameters tend to infinity, the lower argument has
leading coefficient $1/36$ in $m^2n^2$, while the desired upper
expression has coefficient $1/16$. Merely counting forced
$K_{3,3}$ crossings therefore leaves a substantial gap. It allows
different subgraphs to reuse a crossing as efficiently as the
incidence count permits and does not capture global rotation
constraints around vertices.

\paragraph{Verification and final status.}
The root script uses exact integer orientation tests to count
crossing pairs in axial segment drawings for $1\le m,n\le12$
and verifies the displayed small equalities. These tests certify
the construction and arithmetic, not optimality at unproved
parameters. The full matching lower bound remains unresolved
by this attempt.

\subsection{Sources}
\begin{itemize}
\item[S1] ULAM.AI and the UnsolvedMath Contributors, supplied problems.json, record 3319 (OPG-310), OpenGarden collection. Catalogue identity and source transcription; CC BY 4.0.
\url{https://huggingface.co/datasets/ulamai/UnsolvedMath}.
\item[S2] Open Problem Garden, The Crossing Number of the Complete Bipartite Graph. Original problem and discussion; the mathematical formulation below is expanded from the supplied source record.
\url{http://www.openproblemgarden.org/op/the_crossing_number_of_the_complete_bipartite_graph}.

\end{itemize}
\end{document}
